\documentclass[11pt]{article}
\usepackage[T1]{fontenc}
\usepackage[utf8]{inputenc}
\usepackage{lmodern,amsmath,amssymb,amsthm}
\usepackage[margin=1in]{geometry}
\usepackage{microtype,hyperref,enumitem}
\hypersetup{colorlinks=true,linkcolor=blue,citecolor=blue,urlcolor=blue,
pdftitle={Proof of the Linear Hadwiger's conjecture}}
\setlength{\emergencystretch}{2em}
\newtheorem{theorem}{Theorem}
\newtheorem{lemma}[theorem]{Lemma}
\newtheorem{proposition}[theorem]{Proposition}
\newtheorem{corollary}[theorem]{Corollary}
\newcommand{\E}{\mathbb E}
\newcommand{\PP}{\mathbb P}
\newcommand{\pos}[1]{\left(#1\right)_+}
\newtheorem*{quantitativetheorem}{Theorem~\ref{thm:quantitative}}
\newtheorem*{bootstraptheorem}{Theorem~\ref{thm:bootstrap}}
\title{A proof of the Linear Hadwiger's conjecture. }
	\author{ Sergey Norin \thanks{Department of Mathematics and Statistics, McGill University. Supported by an NSERC Discovery grant.}}
\date{}
\begin{document}
\maketitle

\begin{abstract}
We prove that there is an absolute constant $C$ such that every graph
with no $K_t$ minor has chromatic number at most $Ct$. 

We first show that $$\chi(G)\le4h(G)+\varepsilon|G|+A_\varepsilon,$$
where $h(G)$ is the largest order of a complete minor and
$\log A_\varepsilon=O(\varepsilon^{-2}\log(1/\varepsilon))$ as
$\varepsilon\to0$. We then show that a local linear coloring bound
for graphs of order at most $t(\log t)^\alpha$, for any fixed
$\alpha>0$, implies such a bound at exponent $4\alpha/3$.
Iterating this improvement and applying the small-graph
reduction of Delcourt and Postle proves Linear Hadwiger's conjecture.
\end{abstract}

\section{Introduction}

All graphs are finite and simple. A graph $H$ is a minor of a graph
$G$ if a graph isomorphic to $H$ can be obtained from a subgraph of
$G$ by contracting edges. We write $K_t$ for the complete graph on
$t$ vertices. In 1943, Hadwiger~\cite{Had} conjectured that every
$K_t$-minor-free graph is $(t-1)$-colorable. The conjecture holds for
$t\le6$. Wagner's theorem~\cite{Wag} makes the case $t=5$ equivalent
to the Four Color Theorem, and Robertson, Seymour, and
Thomas~\cite{RST} proved the case $t=6$ by a further reduction to
that theorem.

The linear weakening asks whether an absolute constant $C$ suffices
to color every $K_t$-minor-free graph with at most $Ct$ colors.
This question appears explicitly in Reed and Seymour~\cite[p.~148]{RS}
and is formulated as \emph{Linear Hadwiger's conjecture} in
Norin, Postle, and Song~\cite[Conjecture~1.2]{NPS} and Delcourt and
Postle~\cite[Conjecture~1.2]{DP}. Reed and Seymour proved a fractional
version: the fractional chromatic number of every $K_t$-minor-free
graph is at most $2(t-1)$~\cite[Theorem~1.3]{RS}.

For ordinary coloring, Kostochka~\cite{Kos} and Thomason~\cite{Tho}
independently obtained the bound $O(t\sqrt{\log t})$ in the 1980s.
Their results bound the average degree of every $K_t$-minor-free
graph, and hence give upper bounds of this order for its degeneracy
and chromatic number. The density bound is sharp up to a constant
factor, so a better coloring bound requires more than a stronger
bound on degeneracy. Norin, Postle, and Song~\cite[Theorem~1.4]{NPS}
broke this barrier by proving $O(t(\log t)^{1/4+\delta})$-colorability
for every fixed $\delta>0$.

Delcourt and Postle~\cite[Theorem~1.5]{DP} subsequently obtained
$O(t\log\log t)$-colorability. They also reduced the linear conjecture
to graphs of order $O(t(\log t)^4)$ and proved it for graphs excluding
any fixed complete graph as a subgraph
\cite[Theorem~1.6 and Corollaries~1.7--1.10]{DP}.
More recently, Liu and Luo~\cite[Theorem~1.1]{LL} improved the general
bound to $O(t\log\log\log t)$. Lin~\cite[Theorem~1.4 and p.~3]{Lin}
improved the bound on the order of the highly connected subgraphs
used in the Delcourt--Postle reduction, lowering its order cutoff to
$O(t(\log t)^2)$. We prove the linear conjecture.

\begin{theorem}\label{thm:linear}
There is an absolute constant $C$ such that, for every integer $t\ge2$,
every $K_t$-minor-free graph $G$ satisfies $\chi(G)\le Ct$.
\end{theorem}

We write $|G|$, $\chi(G)$,
$\alpha(G)$, and $h(G)$ for the order, chromatic number, independence
number, and largest order of a complete minor of $G$, respectively.
These parameters are zero for the empty graph. A stable set is an
independent set. All logarithms are natural.

The proof has two parts. First, we approximate ordinary coloring by
fractional coloring after adding a graph of bounded maximum degree.
A weighted matching contraction controls the resulting fractional
chromatic number in terms of the complete minors of the original
graph. This yields a local linear coloring bound at a positive
logarithmic order exponent. Second, we improve that exponent by
packing connected induced bipartite subgraphs and localizing chromatic
number near a short induced path. An adaptation of the recursive
assembly of Delcourt and Postle completes this improvement. Repeating
it reaches the local range required by their small-graph reduction.
Section~\ref{sec:deduction} states the two auxiliary theorems and
derives Theorem~\ref{thm:linear} from them; their proofs occupy
Sections~\ref{sec:quantitative} and~\ref{sec:bootstrap}.

\subsection*{AI disclosure}

The proof presented in this paper was developed by OpenAI’s ChatGPT model Astra 6, following research directions and feedback supplied by the authors. The authors suggested several possible approaches, but almost none of their specific proof ideas were retained in the final argument. (Their suggestion to further exploit Liu and Luo’s  ideas of contractions involving unusually large independent sets in vertex neighbourhoods is reflected in the packing argument in Section 4.1.)

The project began with the authors asking the model to prove the very dense case of Linear Hadwiger’s conjecture: there exists an absolute constant $C>0$ such that, for every $C'>0$, there exists $t_0>0$ such that, for every integer $t\ge t_0$, every graph $G$ satisfying
$
|V(G)|\le C't$ and
$
\chi(G)>Ct
$
contains a $K_t$ minor. After the model succeeded, the authors asked it to make the dependence of $t_0$ on $C'$ explicit.
This request yielded a proof of linear Hadwiger's conjecture for graphs $G$ with $h(G)=t$ and $|V(G)| \leq t( \log t )^{1/2-o(1)}$, while by Delcourt-Postle result establishing it for graphs with $|V(G)| \leq t (\log t)^{4+o(1)}$ suffices. 

Several approaches to  were tried at this point. We asked the model, in particular,  to find an original idea to bridge the gap between the two bounds, rather than trying to tighten them. This led to the bootstrap argument used in the present paper.


\subsection{Provenance of the proof ideas}

The hypergraph matching argument in Section~\ref{sec:quantitative} follows the
R\"odl nibble method~\cite{Rodl} and Pippenger's almost-perfect-matching
theorem for nearly regular hypergraphs with small codegrees; see
\cite[Theorem~1.1]{EGJ} for the latter statement. We give an explicit
proof because the dependence on the uniformity is needed when
choosing $\varepsilon$. Independent vertex deletions that equalize
survival probabilities occur in Kang, K\"uhn, Methuku, and
Osthus~\cite[Section~5.1]{KKMO} and Gould and
Kelly~\cite[Section~6.1]{GK}. Our implementation uses Poisson marks
and controls the total degree deficit after trimming. We apply the
resulting matching estimate to an auxiliary hypergraph whose edges
encode stable sets and whose additional vertices control the number
of colors.

The passage from this rounding estimate to the quantitative coloring
bound uses two further constructions proved here. A linear program
with constraints on the weights of pairs yields a bounded-degree
augmentation whose fractional chromatic number approximates the
original chromatic number. A weighted matching contraction then
transfers fractional-coloring weights to a minor of the original
graph. Reed and Seymour's theorem~\cite[Theorem~1.3]{RS}, applied
to that minor, supplies the complete-minor bound. Thus their
fractional theorem is a direct input, while the augmentation and
matching contraction provide the additional comparison with ordinary
coloring.

The factor-two lifting of a quotient coloring through connected
bipartite pieces already appears in Reed and Seymour's discussion of
chordal decompositions~\cite[p.~149]{RS}. Liu and
Luo~\cite[Lemma~3.1]{LL} use a related contraction of induced stars
with independent centers, for which the coloring loss is only one
additional color. In Section~\ref{sec:bootstrap}, a maximal packing
of connected induced bipartite subgraphs of prescribed order leaves
a remainder with no large such subgraph. We then adapt the induced-path
construction of Gy\'arf\'as~\cite[proof of Theorem~2.4]{Gyarfas}:
following a component of large chromatic number localizes the
chromatic obstruction in the neighborhood of a short path. The
absence of large induced stars bounds the independence number of
that neighborhood, and the Reed--Seymour bound converts this into
the order bound needed for the local coloring hypothesis.

The outer recursion is a direct adaptation of Delcourt and
Postle~\cite[Definition~2.4 and Section~7]{DP}. We use their chromatic
separability framework and their assembly of a complete minor from
three smaller models, together with their notion of wovenness and
linkage rerouting lemma~\cite[Definition~5.12 and Lemma~5.15]{DP}.
Our fan lemma follows from their redundant form of Menger's
theorem~\cite[Lemma~5.3]{DP}. The routing hub is obtained from the
connectivity theorem of Gir\~ao and Narayanan~\cite{GN}, in the
quantitative form recorded in~\cite[Theorem~5.1]{DP}; Nguyen's
stronger theorem~\cite[Theorem~1.2]{N} also suffices. Choosing the
hub inside a subgraph of prescribed chromatic number controls the
chromatic cost of its deletion. This permits the recursion to use
the separation supplied by the packing and path argument. The
remaining linkage and rooted-minor inputs are the Thomas--Wollan
theorem~\cite[Corollary~1.3]{TW} and Kawarabayashi's lemma in the
form stated in~\cite[Lemma~5.14]{DP}.

\section{Two auxiliary theorems and the main deduction}\label{sec:deduction}

Our first auxiliary theorem supplies the initial local coloring bound.
\begin{theorem}\label{thm:quantitative}
For every graph $G$ and every $0<\varepsilon\le1$,
\begin{equation}\label{eq:quantitative}
 \chi(G)\le4h(G)+\varepsilon|G|+A_\varepsilon,
 \qquad A_\varepsilon=(100/\varepsilon)^{2000/\varepsilon^2}.
\end{equation}
\end{theorem}

For $\alpha>0$, say that a \emph{local linear bound holds at exponent
$\alpha$} if there are constants $D\ge1$ and $t_0$ such that, for every
integer $t\ge t_0$,
\[
 G\text{ is }K_t\text{-minor-free},\quad |G|\le t(\log t)^\alpha
 \quad\Longrightarrow\quad \chi(G)\le Dt.
\]
The constants may depend on the fixed exponent. The next theorem
improves any such bound to a larger exponent.

\begin{theorem}\label{thm:bootstrap}
If a local linear bound holds at exponent $\alpha>0$, then a local
linear bound holds at exponent $4\alpha/3$.
\end{theorem}

We prove Theorems~\ref{thm:quantitative} and~\ref{thm:bootstrap} in
Sections~\ref{sec:quantitative} and~\ref{sec:bootstrap}, respectively.
We first derive Theorem~\ref{thm:linear} from them. The external input
for this deduction is the following small-graph reduction.

\begin{theorem}[Delcourt--Postle, {\cite[Theorem~1.6]{DP}}]\label{thm:dp}
There is an absolute integer $C_{\rm DP}\ge1$ such that, for every integer
$t\ge3$ and every $K_t$-minor-free graph $G$,
\[
 \chi(G)\le C_{\rm DP}t(1+f(G,t)),
\]
where
\[
 f(G,t)=\max\left(\{0\}\cup\left\{\frac{\chi(H)}a:
 \begin{array}{l}
 H\subseteq G\text{ is }K_a\text{-minor-free},\\
 a\in\mathbb Z,\quad t/\sqrt{\log t}\le a\le t,\\
 |H|\le C_{\rm DP}a(\log a)^4
 \end{array}\right\}\right).
\]
\end{theorem}

We will also use the following explicit complete-minor density bound.
\begin{theorem}[Kostochka--Thomason, {\cite[Theorem~4.3]{DP}}]
\label{thm:density}
Let $r\ge2$ be an integer. Every nonempty graph $H$ with
$|E(H)|/|H|\ge30r\sqrt{\log r}$ contains a $K_r$ minor.
\end{theorem}
Consequently every $K_r$-minor-free graph satisfies
\begin{equation}\label{eq:ktcolor}
 \chi(H)\le60r\sqrt{\log r}+1.
\end{equation}
Indeed, the bound is immediate for the empty graph. Otherwise, an
$s$-chromatic graph has a vertex-critical subgraph $C$ with
minimum degree at least $s-1$, whence
\[
 \frac{|E(C)|}{|C|}\ge\frac{s-1}{2}.
\]
Theorem~\ref{thm:density} applies if the displayed coloring bound fails.

\begin{proof}[Proof of Theorem~\ref{thm:linear}]
Put $\varepsilon=1/[2(\log t)^{1/3}]$. For all sufficiently large $t$,
this lies in $(0,1]$, and
\[
 \log A_\varepsilon
 =\frac{2000}{\varepsilon^2}\log\frac{100}{\varepsilon}
 =8000(\log t)^{2/3}\log\bigl(200(\log t)^{1/3}\bigr)
 =o(\log t).
\]
Consequently $A_\varepsilon\le t/2$ eventually. If $G$ is $K_t$-minor-free
and $|G|\le t(\log t)^{1/3}$, then $h(G)\le t-1$ and
$\varepsilon|G|\le t/2$. Theorem~\ref{thm:quantitative} gives
$\chi(G)\le4(t-1)+t/2+t/2<5t$.

Thus a local linear bound holds at exponent $1/3$. Applying
Theorem~\ref{thm:bootstrap} nine times, each time multiplying the
exponent by $4/3$, gives a local linear bound at exponent
$\alpha=(4/3)^9/3>4$.
Let $D,t_0$ be the constants in this bound. For all sufficiently
large $a$,
\[
 C_{\rm DP}(\log a)^4\le(\log a)^\alpha.
\]
For sufficiently large $t$, every integer $a$ in the window of
Theorem~\ref{thm:dp} satisfies this inequality and $a\ge t_0$.
Thus $f(G,t)\le D$ and $\chi(G)\le C_{\rm DP}(1+D)t$.
The finitely many smaller values of $t$ are absorbed into the constant
using~\eqref{eq:ktcolor}. The case $t=2$
is immediate.
\end{proof}

\section{Proof of the quantitative coloring bound}\label{sec:quantitative}

We prove Theorem~\ref{thm:quantitative} by combining three ingredients.
An explicit hypergraph-matching estimate rounds a fractional coloring
whose pair loads are small. Linear-programming duality then produces
a bounded-degree augmentation whose fractional chromatic number
approximates $\chi(G)$. Finally, a weighted matching contraction
bounds that fractional chromatic number by $4h(G)$, up to an additive
term depending only on the independence number and the augmentation
degree.

A fractional coloring assigns nonnegative weights $x_S$ to the
nonempty stable sets of a graph, with
\[
\sum_{S\ni v}x_S\ge1\qquad(v\in V(G)).
\]
Its cost is $\sum_Sx_S$, and the minimum cost is the fractional
chromatic number $\chi_f(G)$. Linear-programming duality gives
\begin{equation}\label{eq:dual}
\chi_f(G)=\max\left\{\sum_vy_v:
 y_v\ge0,\quad \sum_{v\in S}y_v\le1
 \text{ for every stable set }S\right\}.
\end{equation}
In particular, every dual-feasible weight satisfies $y_v\le1$.
Reed and Seymour~\cite[Theorem~1.3]{RS} proved that every graph $J$
satisfies
\begin{equation}\label{eq:RS}
\chi_f(J)\le2h(J).
\end{equation}
The empty graph satisfies this inequality directly.

\subsection{An explicit almost-perfect matching lemma}

A multihypergraph is $r$-uniform if every edge is a set of $r$ vertices; repeated copies
 of an edge are allowed. The degree of a vertex is the number of incident edge copies,
 and the codegree of two distinct vertices is the number of edge copies containing both.
 A matching consists of pairwise disjoint edges.

\begin{lemma}\label{lem:matching}
Let $r\ge2$ be an integer, let $0<\xi\le1$, and put
\[
 \mu=\left(\frac{\xi}{100r}\right)^{20r}.
\]
Every finite $r$-uniform multihypergraph with positive maximum degree $D_0$, minimum
 degree at least $(1-\mu)D_0$, and maximum codegree at most $\mu D_0$ has a matching
 covering at least a $(1-\xi)$-fraction of its vertices.
\end{lemma}

\begin{proof}
We equalize vertex survival probabilities and control the total degree
deficit after each round. Write $x_+=\max\{x,0\}$.
Let $H_0$ be such a hypergraph, let $N=|V(H_0)|$, and let $B$ be its maximum codegree.
 Since $H_0$ has an edge and $r\ge2$, we have $B\ge1$ and hence $D_0\ge1/\mu$.

Set
\[
 L=\log(4/\xi),\qquad a=\frac{\xi}{100rL},\qquad q=e^{-a},
 \qquad T=\lceil L/a\rceil,
 \qquad D_{i+1}=\lfloor q^{r-1}D_i\rfloor.
\]
We first record bounds ensuring that all $T$ rounds are well defined. We have $1\le
 L\le4/\xi$, $ar\le1/100$, and $aT\le2L$. Put
\[
 b=2\mu(4/\xi)^{2r}.
\]
The definition of $\mu$ gives
\begin{equation}\label{eq:b}
 \sqrt b=\frac{\sqrt2\,4^r\xi^{9r}}{(100r)^{10r}}
 \le \frac{\xi^5}{10^8r^3}.
\end{equation}
Indeed, $9r\ge5$ and
$\sqrt2\,4^r/(100r)^{10r}\le1/(10^8r^3)$ for $r\ge2$.
Since $L\le4/\xi$, we also have $a\ge\xi^2/(400r)$ and hence
\[
 b\le\frac{\xi^{10}}{10^{16}r^6}
   \le\frac{\xi^2}{1600r}\le\frac a4.
\]
The definition of $\mu$ also gives $\mu\le\xi^2/800$.

The geometric sum of the rounding errors in $D_i$ is at most
$1/(1-q^{r-1})\le2/a$, since $r\ge2$ and
$1-e^{-a}\ge a/2$. Thus, for $0\le i\le T$,
\begin{align*}
 D_i&\ge D_0e^{-(r-1)ai}-2/a\\
 &\ge D_0(\xi/4)^{2r}-2/a
 \ge \tfrac12D_0(\xi/4)^{2r}\ge4/a.
\end{align*}
Here $D_0(\xi/4)^{2r}\ge2/b\ge8/a$. Consequently
\begin{equation}\label{eq:degreebounds}
 B/D_i\le b,\qquad D_{i+1}\ge D_i/2,
 \qquad D_{i+1}\ge q^rD_i.
\end{equation}
For the last inequality, note that
\[
 q^{r-1}D_i(1-q)\ge e^{-1/100}(4/a)(a/2)>1,
\]
so $\lfloor q^{r-1}D_i\rfloor\ge q^rD_i$. These bounds also imply $ar^2B/D_i\le1/2$.

We construct a sequence $H_i$ of residual hypergraphs with maximum degree at most $D_i$
 and codegree at most $B$. Write
\[
 n_i=|V(H_i)|,\qquad S_i=n_iD_i-r|E(H_i)|.
\]
The nonnegative quantity $S_i$ is the total degree deficit. Conditional on $H_i$, attach
 to every edge an independent Poisson variable of mean $a/D_i$. Call the edge marked
 when this variable is positive. Independently attach to each vertex $v$ a private
 Poisson variable of mean
\[
 a\left(1-\frac{d_{H_i}(v)}{D_i}\right).
\]
Keep precisely the vertices whose private variable and all incident edge variables are
 zero. Add to the matching every marked edge disjoint from all other marked edges. All
 endpoints of these matching edges are deleted, so the matchings obtained in different
 rounds are disjoint. Private variables do not affect acceptance of an isolated marked
 edge.

Let $J_i$ be induced by the surviving vertices. Delete edges of $J_i$ until its maximum
 degree is at most $D_{i+1}$, and call the result $H_{i+1}$. This trimming does not
 delete vertices or increase codegrees.

Every vertex survives with probability exactly $q$, so
\begin{equation}\label{eq:survivors}
 \E[n_{i+1}\mid H_i]=qn_i,\qquad \E n_i=q^iN.
\end{equation}
More generally, if $U$ is a set of $k$ vertices, independence of the Poisson variables gives
\begin{equation}\label{eq:joint}
 \PP(U\text{ survives}\mid H_i)
 =q^k\exp\!\left(\frac a{D_i}
  \sum_{e\in E(H_i)}\pos{|e\cap U|-1}\right).
\end{equation}
The sum in the exponent is between zero and $\binom{k}{2}B$, since $(j-1)_+\le\binom
 j2$. In particular, every edge survives with probability at least $q^r$.

We next bound the expected number of trimmed edges. Fix $v\in V(H_i)$ and condition on
 its survival. Let $X_v=d_{J_i}(v)$ under this conditioning. Dividing~\eqref{eq:joint}
 for an edge containing $v$ by $q$ yields
\[
 \E X_v\le d_{H_i}(v)q^{r-1}+2ar^2B.
\]
We used $e^x\le1+2x$ for $0\le x\le1/2$ and $d_{H_i}(v)\le D_i$.

For completeness, we bound the conditional variance directly. In the expansion of
 $\operatorname{Var}X_v$, the number of ordered pairs of incident edges sharing a vertex
 other than $v$, including diagonal pairs, is at most
\[
 \sum_{u\ne v}d_{H_i}(u,v)^2
 \le B\sum_{u\ne v}d_{H_i}(u,v)
 \le rBD_i.
\]
Their individual covariance is at most one. For any other pair $e,f$, the sets $e-v$ and
 $f-v$ are disjoint. Conditional on $v$ surviving, its incident edge variables and its
 private variable are forced to be zero; all other variables remain independent. The
 covariance of the survival indicators of these two sets is
\[
 P_eP_f(e^\lambda-1),
\]
where $P_e,P_f$ are their conditional survival probabilities and $\lambda$ is the sum of
 the means of the remaining Poisson variables affecting both sets. Private variables
 cannot affect both. The codegree bound gives $\lambda\le ar^2B/D_i$, and hence the
 covariance is at most $2ar^2B/D_i$. There are at most $D_i^2$ such ordered pairs.
 Therefore
\[
 \operatorname{Var}X_v\le rBD_i+2ar^2BD_i\le4r^2BD_i.
\]
It follows from Cauchy--Schwarz and the rounding in $D_{i+1}$ that
\[
 \E\pos{X_v-D_{i+1}}
 \le2r\sqrt{BD_i}+2ar^2B+1.
\]
All expectations and variances in this paragraph are conditional on $H_i$ and survival of $v$.

Let $F_i$ be the number of edges deleted in trimming. At each deletion, choose an edge
 incident with a vertex whose degree exceeds $D_{i+1}$. Each deletion reduces the sum of
 positive excess degrees by at least one. Thus
\[
 \E[F_i\mid H_i]\le qn_i(2r\sqrt{BD_i}+2ar^2B+1).
\]
Since $qD_{i+1}\le q^rD_i$ and every edge survives with probability at least $q^r$, we obtain
\begin{equation}\label{eq:deficit}
 \E[S_{i+1}\mid H_i]
 \le q^rS_i+rq n_i(2r\sqrt{BD_i}+2ar^2B+1).
\end{equation}

Let $W_i$ be the number of vertices deleted in round $i$ that are not covered by its
 accepted edges. Every such vertex either has a positive private variable or belongs to
 a marked edge intersecting another marked edge. The expected first count is at most
 $aS_i/D_i$. There are at most $|E(H_i)|rD_i$ ordered pairs of distinct intersecting
 edges; each pair is marked with probability at most $(a/D_i)^2$. Charging at most $r$
 vertices to its first edge and using $r|E(H_i)|\le n_iD_i$ gives
\begin{equation}\label{eq:waste}
 \E[W_i\mid H_i]\le aS_i/D_i+ra^2n_i.
\end{equation}

It remains to sum these losses. Put $Z_i=\E S_i/D_i$. By~\eqref{eq:degreebounds}
 and~\eqref{eq:deficit}, using $B\ge1$, $b\le1$, and $ar\le1$,
\[
 \begin{aligned}
 Z_{i+1}
 &\le Z_i+
  \frac{2rN}{D_i}\bigl(2r\sqrt{BD_i}+2ar^2B+1\bigr)\\
 &\le Z_i+
  N\bigl(4r^2\sqrt b+4ar^3b+2rb\bigr)\\
 &\le Z_i+10r^2N\sqrt b.
 \end{aligned}
\]
Here $D_{i+1}\ge D_i/2$ and $\E n_i\le N$ give the first
inequality, while $B/D_i\le b$ and $1/D_i\le B/D_i$ give the second.
Initially $Z_0\le\mu N$, so $Z_i\le\mu N+10ir^2N\sqrt b$. Summing~\eqref{eq:waste}, and
 using $T\le2L/a$, gives
\begin{align*}
 \frac1N\E\sum_{i<T}W_i
 &\le aT\mu+10aT^2r^2\sqrt b+ra^2T\\
 &\le2L\mu+\frac{4000r^3L^3}{\xi}\sqrt b+\frac{\xi}{50}
 <\frac{\xi}{25}.
\end{align*}
For the last inequality, the first term is at most $\xi/100$, and~\eqref{eq:b} and
 $L\le4/\xi$ bound the second by $256000\xi/10^8<\xi/100$. Finally~\eqref{eq:survivors}
 gives $\E n_T=q^TN\le\xi N/4$. The vertices not covered by the union of the accepted
 matchings are precisely those counted by $n_T+\sum_{i<T}W_i$. Their expected number is
 less than $\xi N$, so some outcome proves the lemma.
\end{proof}

\subsection{Fractional coloring and low-codegree rounding}

A fractional coloring can be rounded efficiently when no pair of vertices receives much
 common weight. The color tokens in the construction below control the number of stable
 sets in the resulting coloring.

\begin{lemma}\label{lem:rounding}
Let $m\ge1$ be an integer and let $0<\gamma\le1$. Set
\[
 r=m+2,\qquad
 \mu=\left(\frac{\gamma}{400r^2}\right)^{20r},\qquad
 n_0=\left\lceil10^6r^4\mu^{-4}+2/\gamma\right\rceil.
\]
Suppose $G$ has $n\ge n_0$ vertices, $\alpha(G)\le m$, and a fractional coloring of cost
 $\tau$ satisfying
\[
 \sum_{S\ni v}x_S=1,\qquad
 \sum_{S\supseteq\{u,v\}}x_S\le\mu/4\quad(u\ne v).
\]
Then $\chi(G)\le\tau+\gamma n$.
\end{lemma}

\begin{proof}
Exact vertex loads give
\[
 n=\sum_S |S|x_S,\qquad n/m\le\tau\le n,\qquad x_S\le1.
\]
Take three disjoint vertex sets: $V(G)$, a set $P$ of $p=\lceil\tau\rceil$ color tokens,
 and a set $Z$ of
$z=\lceil(r-1)\tau-n\rceil$ dummy vertices. For every nonempty stable set $S$, token
 $c\in P$, and set $T\subseteq Z$ of size $r-1-|S|$, assign the $r$-set $S\cup\{c\}\cup
 T$ weight
\[
 \omega(S,c,T)=\frac{x_S}{p\binom z{r-1-|S|}}.
\]
Since $1\le |S|\le m$, each edge contains between one and $m$ dummy vertices. Also
\[
 (r-1)\tau-n\ge\tau\ge n/m,
 \qquad n+p+z\le r\tau+2\le rn+2.
\]
Thus $p,z\ge n/m$; the assumed lower bound on $n$ also gives $z\ge m+1$.

The weighted vertex degrees are $1$ at original vertices, $\tau/p$ at tokens, and
 $((r-1)\tau-n)/z$ at dummy vertices. Both $\tau$ and $(r-1)\tau-n$
 are at least $n/m$, while each ceiling adds less than one to its denominator.
 Hence each of these degrees is at least
 \[
 \frac{n/m}{n/m+1}\ge 1-\frac{m}{n},
 \]
 and all are at most one. The weighted codegree
 of two original vertices is at most $\mu/4$. For an original vertex and a token it is
 $1/p$; for an original vertex and a dummy it is at most $m/z$; for a token and a dummy
 it is $((r-1)\tau-n)/(pz)$; and for two dummies it is at most
\[
 \frac{m(m-1)\tau}{z(z-1)}.
\]
Two tokens have codegree zero. For the dummy--dummy bound, for example,
$z\ge n/m$, $z-1\ge n/(2m)$, and $\tau\le n$ give
\[
 \frac{m(m-1)\tau}{z(z-1)}\le\frac{2m^4}{n}.
\]
The other displayed codegrees are no larger, so all weighted codegrees other than
 original--original codegrees are at most $2m^4/n$. Because
$\binom zj\ge z$ for $1\le j\le m$ under our order bound,
\[
 \omega(S,c,T)\le\frac1{pz}\le\frac{m^2}{n^2}.
\]

Independently sample each edge with probability $n\omega(S,c,T)$, which is at most one.
 Chernoff bounds show that, with positive probability, all degrees differ from their
 means by at most $\mu n/8$ and all codegrees are at most $\mu n/2$. Indeed,
$n\ge16m^4/\mu$ gives $2m^4\le\mu n/8$, so every
codegree has mean at most $\mu n/4$. Since there are at most $rn+2$
 vertices, the total failure probability is at most
\[
 8r^2n^2\exp(-\mu^2n/192)<1.
\]
For the strict inequality, $e^x\ge x^4/24$ and
$n\ge10^6r^4\mu^{-4}$ give
\[
 8r^2n^2e^{-\mu^2n/192}
 \le\frac{192^5r^2}{\mu^8n^2}
 \le\frac{192^5}{10^{12}r^6}<1.
\]
Also $n\ge8m/\mu$, so the minimum degree is at least
$n-m-\mu n/8\ge(1-\mu/4)n$, while the maximum degree is at most
$(1+\mu/8)n$. In particular,
\[
 1-\mu/4\ge(1-\mu)(1+\mu/8),
 \qquad \mu/2<\mu(1-\mu/4),
\]
and hence the sampled hypergraph meets both degree hypotheses of
Lemma~\ref{lem:matching}.

Apply Lemma~\ref{lem:matching} with $\xi=\gamma/(4r)$; its parameter is exactly
 $\mu=(\xi/(100r))^{20r}$. The matching leaves at most $\xi(rn+2)\le\gamma n/2$ vertices
 uncovered. It uses at most $p$ edges because each edge contains a distinct token. The
 original vertices on each matching edge form a stable set. Coloring the remaining
 original vertices separately uses at most $p+\gamma n/2\le\tau+\gamma n$ colors, since
 $n\ge2/\gamma$.
\end{proof}

\subsection{A bounded-degree augmentation}

We use the pair-load constraints in Lemma~\ref{lem:rounding} to identify a
 bounded-degree set of nonedges whose addition makes fractional coloring approximate
 ordinary coloring.

\begin{lemma}\label{lem:augmentation}
Let $m\ge2$ be an integer and let $0<\gamma\le1$. Let $\mu,n_0$
 be the parameters in Lemma~\ref{lem:rounding} for $m,\gamma$. Every graph $G$ on $n\ge
 n_0$ vertices with $\alpha(G)\le m$ admits a graph $F$ on $V(G)$ such that
\[
 E(F)\subseteq E(\overline G),\qquad
 \Delta(F)\le\left\lceil\frac{8\binom m2}{\mu\gamma^2}\right\rceil,
\]
and
\[
 \chi_f(G+F)\ge\chi(G)-3\gamma n.
\]
Here $G+F$ is the graph with edge set $E(G)\cup E(F)$.
\end{lemma}

\begin{proof}
Consider the linear program \begin{equation}\label{eq:pairlp}
\begin{aligned}
\tau=\min\;&\sum_S x_S,\\
\text{subject to }&x_S\ge0,\\
&\sum_{S\ni v}x_S\ge1 &&(v\in V(G)),\\
&\sum_{S\supseteq\{u,v\}}x_S\le\mu/4 &&(u\ne v),
\end{aligned}
\end{equation} where \(S\) ranges over nonempty stable sets. Singleton sets give a
 feasible solution of cost \(n\), so \(\tau\le n\).

An optimal solution can be chosen with every vertex load exactly one. Indeed, if a
 vertex is overcovered, transfer appropriate weight from sets containing it to those
 sets with that vertex removed; delete the transferred weight when the resulting set is
 empty. This leaves all other vertex loads unchanged and never increases cost or a pair
 load. Since the original solution was optimal, the resulting feasible solution still
 has cost \(\tau\). Processing all vertices proves the claim. Lemma~\ref{lem:rounding}
 therefore gives \[
\tau\ge\chi(G)-\gamma n.
\]

The dual of \eqref{eq:pairlp} is \begin{equation}\label{eq:pairdual}
\begin{aligned}
\tau=\max\;&\sum_v a_v-\frac{\mu}{4}\sum_{uv}b_{uv},\\
\text{subject to }&a_v,b_{uv}\ge0,\\
&\sum_{v\in S}a_v-\sum_{uv\subseteq S}b_{uv}\le1
&&\text{for every nonempty stable set }S.
\end{aligned}
\end{equation} Only nonedge pairs need to be included. Fix an optimal dual solution.
 Singleton constraints give \(a_v\le1\), and \(\tau\ge0\) gives
 \begin{equation}\label{eq:penalty}
\frac{\mu}{4}\sum_{uv}b_{uv}
 =\sum_v a_v-\tau\le n,
 \qquad\text{so}\qquad
\sum_{uv}b_{uv}\le4n/\mu.
\end{equation} Let \[
X=\left\{v:\sum_u b_{uv}>\frac8{\mu\gamma}\right\}.
\] By \eqref{eq:penalty},
\[
 \frac8{\mu\gamma}|X|
 <\sum_{v\in X}\sum_u b_{uv}
 \le2\sum_{uv}b_{uv}\le\frac{8n}{\mu},
\]
so \(|X|\le\gamma n\). Define \(F\) by adding every nonedge \(uv\)
 with \(u,v\notin X\) and \[
b_{uv}\ge\frac{\gamma}{\binom m2}.
\] Each vertex outside \(X\) has total incident penalty at most \(8/(\mu\gamma)\), while vertices of
 \(X\) receive no added edges. Consequently \[
\Delta(F)\le\frac{8\binom m2}{\mu\gamma^2},
\]
which is independent of \(n\).

Put \(H=G+F\), and assign weights \[
a'_v=
\begin{cases}
a_v/(1+\gamma),&v\notin X,\\
0,&v\in X.
\end{cases}
\] If \(S\) is stable in \(H\), then \(T=S\setminus X\) is stable in \(G\), has at most
 \(m\) vertices, and every pair in \(T\) has penalty less than
 \(\gamma/\binom m2\). The
 constraint in \eqref{eq:pairdual} therefore yields \[
\sum_{v\in T}a_v\le1+\sum_{uv\subseteq T}b_{uv}
 \le1+\binom{|T|}{2}\frac{\gamma}{\binom m2}
 \le1+\gamma.
\] The same conclusion is immediate if \(T\) is empty. Hence \(a'\) is dual-feasible for
 fractional coloring of \(H\). Using \(\sum_v a_v\ge\tau\), \(|X|\le\gamma n\), and
 \(\tau\le n\), we obtain \[
\begin{aligned}
\chi_f(H)
&\ge\frac{\sum_v a_v-|X|}{1+\gamma}\\
&\ge\frac{\tau-\gamma n}{1+\gamma}\\
&\ge\tau-2\gamma n\\
&\ge\chi(G)-3\gamma n
.
\end{aligned}
\]
For the penultimate inequality, use
$(\tau-\gamma n)/(1+\gamma)
=\tau-\gamma(\tau+n)/(1+\gamma)\ge\tau-2\gamma n$,
since $\tau\le n$. This proves the lemma. \end{proof}

\subsection{Robust fractional bound}

The augmentation changes the graph whose fractional chromatic number is measured. The
 next proposition controls this change using the minor number of the original graph. All
 dependence on the independence number and the degree of the augmentation enters an
 additive constant.

\begin{proposition}[Bounded-degree perturbations]\label{prop:robust} Let \(m\ge1\) and
 \(d\ge0\) be integers. Suppose that \(\alpha(G)\le m\) and that \(F\) is a graph on
 \(V(G)\) with \[
E(F)\subseteq E(\overline G),\qquad \Delta(F)\le d.
\] Set \[
K=\sum_{j=0}^{2m-1}(d+1)^j.
\] Then \begin{equation}\label{eq:robust}
\chi_f(G+F)\le4h(G)+2m(K+1).
\end{equation} \end{proposition}

\begin{lemma}[Independent transversal]\label{lem:transversal} Let \(P\) be a perfect
 matching on a finite set \(U\), and let \(D\) be a graph on \(U\) edge-disjoint from
 \(P\). If no cycle of \(D+P\) contains an edge of \(P\), then \(D\) has an independent
 set containing exactly one endpoint of every edge of \(P\).

\end{lemma}

\begin{proof} Every edge of \(P\) is a bridge of \(D+P\). Contract the connected
 components of \(D\). No edge of \(P\) has both endpoints in one component, since that
 edge together with a path inside the component would form a cycle. The resulting
 multigraph is a forest: a cycle, including two parallel edges, would make its matching
 edges nonbridges in \(D+P\).

The degree of the vertex corresponding to a component \(C\) of \(D\) is exactly \(|C|\).
 Indeed, each vertex of \(C\) is incident with one matching edge and every matching edge
 leaves \(C\). If \(U\ne\varnothing\), this forest has a leaf. The corresponding
 component of \(D\) is therefore a singleton \(\{x\}\), so \(x\) has no neighbor in
 \(D\).

Select \(x\), and delete \(x\) and its matching partner. On the remaining vertices the
 matching is perfect and the hypothesis about cycles is preserved. Induction supplies an
 independent transversal of the remaining pairs, and adjoining \(x\) proves the lemma.
 The empty set provides the base case. Notice that cycles consisting only of edges of
 \(D\) cause no difficulty. \end{proof}

\begin{lemma}[Matching with bounded weight loss]\label{lem:weighted} Let \(m\ge1\) and
 \(d\ge0\) be integers. Let \(G\) and \(D\) be graphs on the same finite vertex set,
 with \(\alpha(G)\le m\) and \(\Delta(D)\le d\). Let \(w(v)\in[0,1]\) be vertex weights.
 Fix an integer \(r\ge1\) and set \[
B=\sum_{j=0}^{r}(d+1)^j.
\] There is a matching \(M\subseteq E(G)\), with unmatched set \(X\), such that
 \begin{equation}\label{eq:matchingloss}
|X|\le m(B+1),\qquad
\sum_{uv\in M}|w(u)-w(v)|\le m(B+1),
\end{equation} and no cycle of length at most \(r+1\) in \(D+M\) contains an edge of
 \(M\). Moreover, \(E(M)\cap E(D)=\varnothing\).

\end{lemma}

\begin{proof} Maintain a set of unprocessed vertices, initially \(V(G)\), and an
 initially empty matching \(M\). Choose an unprocessed vertex \(v\) of maximum weight.
 Its eligible partners are the other unprocessed vertices \(u\) such that \(uv\in E(G)\)
 and the distance from \(u\) to \(v\) in the current graph \(D+M\) is greater than
 \(r\). Vertices in different components have infinite distance. If an eligible partner
 exists, choose one of maximum weight, insert \(uv\) into \(M\), and remove both
 endpoints from the unprocessed set. Otherwise, put \(v\) into \(X\) and remove it from
 the unprocessed set. Continue until all vertices have been processed. Throughout this
 procedure, the graph used to measure distances retains the full vertex set; putting a
 vertex in \(X\) only removes it from consideration as a future partner.

At every stage the graph \(D+M\) has maximum degree at most \(d+1\), and its
 radius-\(r\) balls have at most \(B\) vertices. Every inserted matching edge is a
 nonedge of \(D\), because its endpoints have distance greater than \(r\ge1\).

Suppose that the final graph has a cycle of length at most \(r+1\) containing a matching
 edge. Consider the matching edge of this cycle inserted last. All remaining edges of
 the cycle already existed when this edge was inserted, and form a path of length at
 most \(r\) between its endpoints. This contradicts the insertion rule.

Order the vertices of \(X\) by the time they were put aside. If \(x\in X\), every later
 vertex of \(X\) that is a \(G\)-neighbor of \(x\) was ineligible when \(x\) was
 processed. All these neighbors were therefore in the radius-\(r\) ball about \(x\) in
 the graph \(D+M\) at that time. Consequently \(x\) has at most \(B\) later neighbors in
 \(G[X]\). Greedy coloring in reverse order gives a proper coloring of \(G[X]\) with at
 most \(B+1\) colors. Each color class has at most \(m\) vertices, since \(\alpha(G)\le
 m\), and hence \[
|X|\le m(B+1).
\]

For the weight estimate, fix a threshold \(a\in[0,1]\). Call an edge of \(M\) a crossing
 edge if one endpoint has weight greater than \(a\) and the other has weight at most
 \(a\). Let \(Y_a\) consist of the higher-weight endpoint of each crossing edge. For
 every crossing edge \(vu\), the vertex selected first by the procedure is its
 higher-weight endpoint \(v\), so \[
w(v)>a\ge w(u).
\] Order \(Y_a\) by the insertion times of its matching edges. When \(v\in Y_a\) was
 processed, every later vertex \(y\in Y_a\) was still unprocessed and had weight greater
 than \(a\), whereas the chosen partner \(u\) had weight at most \(a\). If \(vy\in
 E(G)\), then \(y\) must have been ineligible, since an eligible \(y\) would have been
 preferred to \(u\). Every later \(G\)-neighbor of \(v\) in \(Y_a\) was therefore in the
 current radius-\(r\) ball about \(v\). Thus every vertex of \(G[Y_a]\) has at most
 \(B\) later neighbors. As before, \(G[Y_a]\) is \((B+1)\)-colorable, so \[
|Y_a|\le m(B+1).
\] The matching has exactly \(|Y_a|\) crossing edges.

Integrating over thresholds gives \[
\begin{aligned}
\sum_{uv\in M}|w(u)-w(v)|
&=\int_0^1
\bigl|\{uv\in M:\min(w(u),w(v))\le a<\max(w(u),w(v))\}\bigr|\,da\\
&\le m(B+1).
\end{aligned}
\] This proves \eqref{eq:matchingloss}. \end{proof}

\begin{proof}[Proof of Proposition~\ref{prop:robust}] Put \(H=G+F\). Choose an optimal
 weight function \(w\) in~\eqref{eq:dual} for \(H\). Thus \[
\sum_{v\in V(G)}w(v)=\chi_f(H),\qquad 0\le w(v)\le1.
\] Apply Lemma~\ref{lem:weighted} to \(G\), \(D=F\), these weights, and \(r=2m-1\). Let
 \(M\) be the resulting matching and let \(X\) be its unmatched set. Since \(M\subseteq
 E(G)\), \[
Q=(G-X)/M
\] is a minor of \(G\). Here the quotient retains every edge arising from an edge of
 \(G-X\), removing only loops and identifying parallel edges. Assign to the quotient
 vertex corresponding to an edge \(uv\in M\) the weight \[
z(uv)=\min\{w(u),w(v)\}.
\]

We claim that \(z\) is feasible in~\eqref{eq:dual} for \(Q\). Let \(S\) be any
 independent set of \(Q\), and let \(U\) be the union of its matching pairs. There is no
 \(G\)-edge between different pairs of \(S\), by the definition of this quotient.
 Selecting one endpoint from each pair therefore gives an independent set in \(G\),
 whence \[
|S|\le m,\qquad |U|\le2m.
\] Furthermore, \(G[U]\) consists exactly of the matching \(M[U]\).

No cycle of \(F[U]+M[U]\) contains a matching edge: such a cycle would have length at
 most \(|U|\le2m\), contrary to Lemma~\ref{lem:weighted}. Lemma~\ref{lem:transversal}
 therefore gives a set \(T\subseteq U\) that is independent in \(F\) and contains
 exactly one endpoint of every pair in \(S\). Since \(G[U]\) consists of the matching
 pairs, \(T\) is also independent in \(G\), and hence in \(H\). Therefore \[
\sum_{e\in S}z(e)\le\sum_{v\in T}w(v)\le1.
\] This proves feasibility.

Lemma~\ref{lem:weighted} now yields \[
\begin{aligned}
2\sum_{e\in M}z(e)
&=\sum_{v\notin X}w(v)-\sum_{uv\in M}|w(u)-w(v)|\\
&\ge\chi_f(H)-|X|-m(K+1)\\
&\ge\chi_f(H)-2m(K+1).
\end{aligned}
\] By fractional-coloring duality, \[
\chi_f(H)\le2\chi_f(Q)+2m(K+1).
\] Finally, \eqref{eq:RS} and \(h(Q)\le h(G)\) imply \[
\chi_f(H)\le4h(Q)+2m(K+1)\le4h(G)+2m(K+1),
\] which is \eqref{eq:robust}. \end{proof}

\subsection{Completing the proof}

\begin{quantitativetheorem}
For every graph $G$ and every $0<\varepsilon\le1$,
\[
 \chi(G)\le4h(G)+\varepsilon|G|+(100/\varepsilon)^{2000/\varepsilon^2}.
\]
\end{quantitativetheorem}
\begin{proof}
Fix $0<\varepsilon\le1$ and set
\[
 m=\lceil2/\varepsilon\rceil,\qquad r=m+2,\qquad
 \gamma=\varepsilon/6,\qquad
 \mu=\left(\frac{\varepsilon}{2400r^2}\right)^{20r}.
\]
Lemma~\ref{lem:rounding} gives the pair-load tolerance $\mu/4$ and order threshold
\begin{equation}\label{eq:n0}
 n_0=\left\lceil10^6r^4\mu^{-4}+12/\varepsilon\right\rceil.
\end{equation}
Since $m\ge2$, Lemma~\ref{lem:augmentation} applies. Every graph $G$ on $n\ge n_0$
 vertices with $\alpha(G)\le m$ admits $F\subseteq\overline G$ such that
\[
 \Delta(F)\le d:=\left\lceil\frac{288\binom m2}{\mu\varepsilon^2}\right\rceil,
 \qquad
 \chi(G)\le\chi_f(G+F)+\varepsilon n/2.
\]
Proposition~\ref{prop:robust} now gives
\[
 \chi(G)\le4h(G)+\varepsilon n/2+A,
\]
where
\[
 A=\max\left\{n_0,\,
 2m\left(1+\sum_{j=0}^{2m-1}(d+1)^j\right)\right\}.
\]
The same inequality holds when $n<n_0$, since $\chi(G)\le n<n_0$. Thus it holds for
 every graph of independence number at most $m$, including the empty graph.

For an arbitrary graph $G$ on $n$ vertices, repeatedly remove a stable set of size
 exactly $m+1$ while such a set exists. There are at most
$n/(m+1)$ removed sets; let $J$ be the remaining induced subgraph. Then
\[
 \frac{n}{m+1}\le\frac{\varepsilon n}{2},
 \qquad \alpha(J)\le m,\qquad h(J)\le h(G).
\]
Color the removed sets separately and apply the preceding bound to $J$. This gives
\[
 \chi(G)\le\frac{n}{m+1}+\chi(J)
 \le\frac{\varepsilon n}{2}+4h(J)+\frac{\varepsilon |V(J)|}{2}+A
 \le4h(G)+\varepsilon n+A.
\]

Finally, we bound that fully explicit expression. Put $R=100/\varepsilon$. Since
 $m\le3/\varepsilon$ and $r\le5/\varepsilon$,
\[
 \mu^{-1}
 =\left(\frac{2400r^2}{\varepsilon}\right)^{20r}
 \le(60000\varepsilon^{-3})^{100/\varepsilon}
 \le R^{300/\varepsilon}.
\]
The last step uses $60000\varepsilon^{-3}\le R^3$.
Also, $\binom m2\le9/(2\varepsilon^2)$, so the ceiling in the
definition of $d$ yields
\[
 d+1\le\frac{1298}{\varepsilon^4}\mu^{-1}
 \le R^{304/\varepsilon}.
\]
Here $1298\varepsilon^{-4}\le R^{4/\varepsilon}$.
Using $\sum_{j=0}^{2m-1}(d+1)^j\le2m(d+1)^{2m}$ and $2m\le6/\varepsilon$ gives
\[
 2m\left(1+\sum_{j=0}^{2m-1}(d+1)^j\right)
 \le\left(\frac6\varepsilon+\frac{36}{\varepsilon^2}\right)
           (d+1)^{2m}
 \le42\varepsilon^{-2}R^{1824/\varepsilon^2}
 \le R^{1826/\varepsilon^2}.
\]
The final step uses $42\varepsilon^{-2}\le R^{2/\varepsilon^2}$.
For the other term in $A$, \eqref{eq:n0} and $\varepsilon\le1$ give
\[
 n_0\le10^6(5/\varepsilon)^4R^{1200/\varepsilon}
          +12/\varepsilon+1
 \le\bigl(6.25\cdot10^8\varepsilon^{-4}
           +13\varepsilon^{-1}\bigr)R^{1200/\varepsilon}
 \le R^{1206/\varepsilon^2}.
\]
Therefore $A\le R^{2000/\varepsilon^2}$, which
proves~\eqref{eq:quantitative}.

\end{proof}



\section{Proof of the exponent improvement}\label{sec:bootstrap}

We prove Theorem~\ref{thm:bootstrap}. Packing connected induced
bipartite subgraphs produces a smaller minor, which is colored using
the assumed local bound. In the remainder, a short induced path
localizes chromatic number whenever separation fails. The same local
bound then supplies the separation needed for a recursive minor
construction. The construction follows Delcourt and Postle, with the
chromatic number of its routing subgraph controlling the cost of deletion.

\subsection{Packing and localization}

Let $b(G)$ be the largest order of a connected induced bipartite
subgraph of $G$, with $b(\varnothing)=0$. This parameter is inherited
by induced subgraphs. 

The first reduction separates a union of bipartite subgraphs from a
remainder in which $b$ is small.
\begin{lemma}\label{lem:packing}
For every graph $G$ and integer $k\ge2$, there is a partition
$V(G)=W\mathbin{\dot\cup}R$ and a minor $Q$ of $G$ such that
\[
 |Q|\le\lfloor |G|/k\rfloor,\qquad b(G[R])<k,
 \qquad \chi(G[W])\le2\chi(Q).
\]
\end{lemma}
\begin{proof}
Greedily choose disjoint sets of $k$ vertices inducing connected
bipartite graphs, until no more can be chosen. Every connected graph
of order at least $k$ has a connected induced subgraph of order exactly
$k$: delete leaves of a spanning tree until $k$ vertices remain.
Thus the remaining induced graph has $b<k$.
Let $W$ be the union of the chosen sets. Contract each chosen set in
$G[W]$ to obtain $Q$, retaining all adjacencies between different sets.
A coloring of $Q$, together with a separate bipartition bit inside
each contracted set, is a coloring of $G[W]$ with at most $2\chi(Q)$
colors. The empty packing is allowed.
\end{proof}

When chromatic separation fails, the next lemma localizes the excess
chromatic number in the closed neighborhood of a short induced path.

\begin{lemma}\label{lem:path}
Let $q\ge1$ and $k\ge2$ be integers. Suppose that $b(G)\le k$,
$\chi(G)\ge q$, and $G$ has no two vertex-disjoint subgraphs of chromatic
number at least $q$. Then there is an induced path $P$ of order at most
$k$ such that, with $J=G[N_G[V(P)]]$,
\[
 \chi(G-V(J))\le q-1,\qquad
 \alpha(J)\le k(k-1),\qquad
 \chi(J)\ge\chi(G)-q+1.
\]
Here $N_G[X]$ denotes $X$ together with all vertices having a neighbor in $X$.
\end{lemma}
\begin{proof}
At most one component of any induced subgraph of $G$ has chromatic
number at least $q$, since two such components would violate the
hypothesis. Let $C_0$ be the unique such component of $G$, and choose
$v_1\in V(C_0)$.

We grow an induced path $P_i=v_1\cdots v_i$. If all components of
$G-N_G[V(P_i)]$ have chromatic number less than $q$, we stop.
Otherwise let $C_i$ be its unique component of chromatic number at
least $q$. The construction will ensure that $C_i$ is contained in
$C_{i-1}$ and that $v_i$ has a neighbor in $C_{i-1}$ when $i\ge2$.
For $i=1$, connectivity of $C_0$ supplies a vertex
$v_2\in N_G(v_1)$ adjacent to $C_1$.
For $i\ge2$, the component $C_i$ is a component of
$C_{i-1}-N_G(v_i)$. It is a proper subset, because $v_i$ has a neighbor
in $C_{i-1}$. Since $C_{i-1}$ is connected, there is a vertex
\[
 v_{i+1}\in V(C_{i-1})\cap N_G(v_i)
\]
adjacent to $C_i$. This vertex lies outside the closed neighborhoods
of $v_1,\ldots,v_{i-1}$, so the extended path is induced.
Every component discarded earlier had chromatic number less than
$q$; deleting further vertices cannot increase that number.
Hence any surviving component of chromatic number at least $q$
lies in $C_i$, establishing the invariant at the next step.

An induced path is bipartite, so the process cannot extend to $k+1$
vertices. At termination its closed neighborhood therefore has the
required complement. Also $\alpha(G[N_G[v]])\le k-1$ for every vertex
$v$: a stable set of $k$ neighbors would form an induced star of order
$k+1$, and a stable set in the closed neighborhood containing $v$
has order one. Taking the union of the closed neighborhoods of the
at most $k$ path vertices gives $\alpha(J)\le k(k-1)$.
Finally $\chi(G)\le\chi(J)+\chi(G-V(J))$ gives the last inequality.
\end{proof}

The fractional bound~\eqref{eq:RS} converts the independence bound
into an order bound. If $(x_S)$ is an optimal fractional coloring of
$H$, summing its covering constraints gives
\begin{equation}\label{eq:order}
 |H|\le\sum_S|S|x_S
 \le\alpha(H)\sum_Sx_S
 =\alpha(H)\chi_f(H)
 \le2\alpha(H)h(H).
\end{equation}
This lets us turn a local coloring bound into chromatic separation.

\begin{lemma}\label{lem:separation}
Suppose every $K_u$-minor-free graph of order at most $2uk^2$ has
chromatic number at most $du$, where $u,d,k$ are positive integers and
$k\ge2$. If $Y$ is $K_u$-minor-free, $b(Y)\le k$, and $\chi(Y)>du$,
then $Y$ has two vertex-disjoint induced subgraphs, each of chromatic
number at least $\chi(Y)-du$.
\end{lemma}
\begin{proof}
If the asserted pair does not exist, apply Lemma~\ref{lem:path} with
$q=\chi(Y)-du\ge1$. The resulting induced graph $J$ satisfies
\[
 \alpha(J)\le k(k-1),\qquad
 \chi(J)\ge\chi(Y)-(\chi(Y)-du)+1=du+1.
\]
It is $K_u$-minor-free, so \eqref{eq:order} gives
\[
 |J|\le2\alpha(J)h(J)
 \le2k(k-1)(u-1)<2uk^2.
\]
This contradicts the assumed local bound.
\end{proof}

\subsection{Connectivity tools and a hub of bounded chromatic number}

We state the results needed for the recursive construction. All references
to Delcourt and Postle use the numbering in~\cite{DP}.
Graphs are finite and simple, $|H|=|V(H)|$, and logarithms are natural.
A graph is $r$-connected if it has more than $r$ vertices and remains
connected after deleting fewer than $r$ vertices. We denote its vertex
connectivity by $\kappa(H)$ and its minimum degree by $\delta(H)$.
A $K_r$ model consists of $r$ pairwise disjoint nonempty connected branch
sets, every two of which are joined by an edge. It is rooted at $r$
specified distinct vertices if each branch set contains its corresponding
root. A linkage is a family of pairwise vertex-disjoint paths joining
prescribed pairs. A path from a set $U$ to a disjoint set $V$ meets these
sets only at its corresponding ends. For a path family $\mathcal P$ or
a model $\mathcal M$, write $V(\mathcal P)$ or $V(\mathcal M)$ for the
union of its vertex sets.

The first theorem restores connectivity after vertices have been deleted.
\begin{theorem}[Gir\~ao--Narayanan, {\cite[Theorem~5.1]{DP}}]
\label{thm:connectivity}
Let $r\ge1$ be an integer. Every graph $H$ with $\chi(H)\ge7r$
contains an $r$-connected subgraph $H'$ with
$\chi(H')\ge\chi(H)-6r$.
\end{theorem}
We may take the induced closure of $H'$, since this preserves connectivity
and cannot decrease its chromatic number.

We use Menger's theorem and the following quantitative linkedness result.
\begin{theorem}[Thomas--Wollan, {\cite[Theorem~5.6 and the following remark]{DP}}]
\label{thm:linked}
Let $\ell\ge1$ be an integer. Every $10\ell$-connected graph contains
a linkage joining any prescribed $\ell$ pairs of distinct terminals.
\end{theorem}

A sufficiently large complete minor also supplies prescribed roots.
\begin{lemma}[Kawarabayashi, {\cite[Lemma~5.14]{DP}}]\label{lem:rooted}
Let $r\ge1$ be an integer. Every $r$-connected graph containing a
$K_{2r}$ minor has a $K_r$ model rooted at any $r$ specified distinct
vertices.
\end{lemma}

For the initial scale we use the density theorem and its coloring
consequence~\eqref{eq:ktcolor}, stated before the proof of
Theorem~\ref{thm:linear}.

The next lemma supplies a hub whose deletion has bounded chromatic cost.
\begin{lemma}\label{lem:hub}
Let $r\ge1$ be an integer. If $\chi(H)\ge7r$, then $H$ has an
induced $r$-connected subgraph $H_0$ with $\chi(H_0)\le7r$.
\end{lemma}
\begin{proof}
Delete vertices one at a time to obtain an induced subgraph $X$ with
$\chi(X)=7r$. Apply Theorem~\ref{thm:connectivity} to $X$ and take the
induced closure of the resulting subgraph.
\end{proof}

To join the hub to both the prescribed terminals and the remaining graph,
we use Delcourt and Postle's redundancy lemma.
\begin{lemma}[Delcourt--Postle, {\cite[Lemma~5.3]{DP}}]
\label{lem:redundancy}
Let $A_1,A_2,B$ be disjoint nonempty vertex sets in a graph $H$.
For each $i\in\{1,2\}$, suppose that $H-A_{3-i}$ contains $2|A_i|$
paths from $A_i$ to $B$, pairwise disjoint outside $A_i$, with exactly
two paths starting at each vertex of $A_i$.
Then $H$ contains $|A_1|+|A_2|$ pairwise vertex-disjoint paths from
$A_1\cup A_2$ to $B$.
\end{lemma}

We record the consequence used below.
\begin{lemma}\label{lem:fan}
Let $a\ge1$ be an integer, and let $Z,U,V$ be disjoint vertex sets
in a graph, with $|Z|=7a$. Suppose there are $14a$ paths from $Z$ to
$V$, two from each vertex of $Z$, disjoint outside $Z$ and avoiding
$U$, and $2a$ disjoint paths from $U$ to $V$ avoiding $Z$.
Then there are $8a$ disjoint paths to $V$, one from each vertex of
$Z$ and $a$ from distinct vertices of $U$, with all interiors outside
$Z\cup U\cup V$.
\end{lemma}
\begin{proof}
Work in the subgraph consisting of the given paths and their edges.
Pair the $2a$ starting vertices in $U$ of the second path family.
For each pair add a new vertex adjacent to its two members, and let
$A_2$ be the set of the $a$ new vertices. Extending the second path
family by these edges gives two paths from each vertex of $A_2$ to
$V$, disjoint outside $A_2$ and avoiding $Z$. The first family avoids
$A_2$. Lemma~\ref{lem:redundancy}, with $A_1=Z$ and $B=V$, gives
$8a$ disjoint paths, one from each vertex of $Z\cup A_2$.

Delete the new starting vertices. Vertices of $U$ have degree at most one
in the union of the original paths, so they cannot be internal vertices
of the resulting paths. The other exclusions hold by construction.
\end{proof}

We use the rooted convention of~\cite[Definition~5.12]{DP}.
For integers $a\ge1$ and $b\ge0$, a graph is \emph{$(a,b)$-woven} if,
for any distinct roots $r_1,\ldots,r_a$ and terminal pairs
$(s_1,t_1),\ldots,(s_b,t_b)$, with different pairs having no common
terminal but allowing $s_i=t_i$, it contains a $K_a$ model rooted at
the roots and a linkage for the pairs whose vertex intersection is
exactly the roots that are also terminals. A singleton pair is represented
by a path of length zero. Roots may coincide with terminals.

The following form of the weaving lemma includes the containment
properties established in its cited proof.
\begin{lemma}[Delcourt--Postle, {\cite[Lemma~5.15 and its proof]{DP}}]
\label{lem:weave}
Let $a,b\ge1$ be integers. Let $H$ be an $(a,b)$-woven subgraph of a
graph $F$, with $|H|\ge a+2b$, and let $R\subseteq V(H)$ have size
$a$. If $\mathcal P$ is a linkage in $F$ for $b$ pairs with all $2b$
terminals distinct, then there are a $K_a$ model $\mathcal M$ in
$H$ rooted at $R$ and a linkage $\mathcal P'$ for the same pairs such that
\[
 V(\mathcal P')\subseteq V(\mathcal P)\cup V(H),\qquad
 V(\mathcal P')\cap V(\mathcal M)\subseteq R\cap V(\mathcal P).
\]
\end{lemma}

The following simultaneous form follows from Delcourt and Postle's
linkage rerouting argument~\cite[Claim~7.1.1]{DP}.
\begin{lemma}
\label{lem:simultaneous}
Let $a$ be a positive integer divisible by three, and put $a'=2a/3$.
Suppose that a graph $F$ contains disjoint $(a',2a)$-woven subgraphs
$J_1,J_2,J_3$, each of order at least $a'+4a$.
Let $s_1,\ldots,s_{2a}$ be distinct vertices outside their union, and
let $t_1,\ldots,t_{2a}$ be distinct vertices such that
\[
 T_i=\{t_{(i-1)a'+1},\ldots,t_{ia'}\}\subseteq V(J_i)
 \quad(i=1,2,3).
\]
If $F$ has a linkage for the pairs $(s_j,t_j)$, then it has a linkage
$\mathcal P$ for the same pairs and $K_{a'}$ models $\mathcal M_i$
in $J_i$, rooted at $T_i$, such that
$V(\mathcal P)\cap V(\mathcal M_i)=T_i$ for $i=1,2,3$.
\end{lemma}
\begin{proof}
Starting with the given linkage, apply Lemma~\ref{lem:weave} to
$J_1,J_2,J_3$ in that order, always with all $2a$ terminal pairs.
At each application, the linkage meets the new model only at roots
that were already on the linkage. All $a'$ roots in $T_i$ are
endpoints of its paths, so the intersection is exactly $T_i$.
The new linkage lies in the union of the previous one and $J_i$.
Since the three subgraphs are disjoint, rerouting in $J_i$ cannot
create an intersection with either of the models constructed earlier.
Their roots remain on the linkage as endpoints.
\end{proof}

\subsection{The outer recursion with a separation hypothesis}

For a nonnegative integer $s$, a graph $Y$ is
\emph{$s$-chromatic-separable} if it has two vertex-disjoint subgraphs,
each of chromatic number at least $\chi(Y)-s$
\cite[Definition~2.4]{DP}. We may require these subgraphs to be induced
by taking their induced closures.

The next lemma adapts the proof of~\cite[Theorem~7.1]{DP}.
Here separation is a hypothesis, and Lemma~\ref{lem:hub} bounds the
chromatic cost of the hub. The linkage and weaving steps are supplied
by the results stated above.
\begin{lemma}\label{lem:outer}
Let $T\ge100$ be a power of three and let $d\ge1$ be an integer.
Suppose that a graph $G$ has the following property: for every integer
scale $a=(2/3)^iT>T/\sqrt{\log T}$, where $i\ge0$, every induced
$K_{14a}$-minor-free subgraph $Y$ of $G$ with $\chi(Y)>28da$ is
$14da$-chromatic-separable.
Put
\[
 K=10000,\qquad A=2000,\qquad B=10^6(d+1).
\]
For every integer scale $a=(2/3)^iT$ with $i\ge0$, each induced subgraph $F$ of $G$
with $\kappa(F)\ge Ka$ and $\chi(F)\ge AT+Ba$ is $(a,3a)$-woven.
\end{lemma}
\begin{proof}
We induct on $a$ among the integer scales. We first separate any
coincident roots and terminals. Let $Z_0$ be the set of vertices
occupying these roles, so $|Z_0|\le7a$. Greedily choose a distinct
neighbor outside $Z_0$ for each of the $7a$ roles, including both
roles of a singleton pair; this is possible since $\delta(F)>14a$.
Call these new vertices the proxies, and put $F_0=F-Z_0$. Then
\[
 \kappa(F_0)\ge(K-7)a,\qquad
 \chi(F_0)\ge AT+(B-7)a.
\]
A model and linkage for the distinct proxies extend to the original
vertices by their chosen edges. For a singleton pair, discard its
proxy path and use its original vertex. Since $F_0$ avoids $Z_0$,
this gives exactly the permitted root--terminal intersections.

If $F_0$ contains a $K_{14a}$ minor, Lemma~\ref{lem:rooted} gives
a $K_{7a}$ model rooted at all the proxies, since
$\kappa(F_0)\ge7a$. Retain its $a$ root branches, and join each pair
of terminal branches by an edge and paths inside those branches.
This gives the required model and linkage.
In particular, this applies when $a\le T/\sqrt{\log T}$, because
\[
 840a\sqrt{\log(14a)}+1
 \le840T\sqrt{\frac{\log(14T)}{\log T}}+1
 \le840T\sqrt{1+\frac{\log14}{\log100}}+1
 <2000T
\]
for $T\ge100$, and~\eqref{eq:ktcolor} forces a $K_{14a}$ minor
in $F_0$. We may therefore assume that $a>T/\sqrt{\log T}$ and
that $F_0$ is $K_{14a}$-minor-free.

\medskip\noindent\emph{Linking through a hub.}
Let $Z$ be the set of $7a$ proxies. Apply Lemma~\ref{lem:hub} in
$F_0-Z$ with $r=140a$ to obtain an induced subgraph $H_0$ with
\[
 \kappa(H_0)\ge140a,\qquad \chi(H_0)\le980a.
\]
The hypothesis holds because $\chi(F_0-Z)\ge AT+(B-14)a\ge980a$.
Since $\kappa(F_0)\ge(K-7)a\ge14a$ and $|H_0|\ge14a$,
Menger's theorem,
with each source in $Z$ replaced by two copies, gives $14a$ paths
from $Z$ to $H_0$, two from each source and disjoint outside $Z$.
Choose each path induced by shortcutting it. The graph induced by
their union has chromatic number at most $28a$, using a separate
two-color palette for each path.

Delete $Z$, $H_0$, and these paths from $F_0$, and apply
Theorem~\ref{thm:connectivity} with $r=Ka$ to the remaining induced
graph. Its chromatic number is at least $AT+(B-1022)a>7Ka$.
We obtain an induced subgraph $F_3$ with
\[
 \kappa(F_3)\ge Ka,\qquad
 \chi(F_3)\ge AT+(B-1022-6K)a.
\]

It suffices to prove that $F_3$ is $(a,0)$-woven. To see this, apply
Menger's theorem in $F_0-Z$ to find $2a$ disjoint paths
from $F_3$ to $H_0$; the host is at least $2a$-connected because
$\kappa(F_0-Z)\ge(K-14)a\ge2a$, and both
endpoint sets have size at least $2a$. The earlier double fan avoids
$F_3$, so Lemma~\ref{lem:fan} gives $8a$ disjoint paths to $H_0$,
one from each proxy and $a$ from distinct vertices of $F_3$, with
interiors outside these endpoint sets. Root a $K_a$ model in $F_3$
at the latter ends. By Theorem~\ref{thm:linked}, the hub has a
linkage pairing its $8a$ ends as required: $3a$ pairs connect the
terminal paths, and $a$ pairs connect the root paths to the paths
from $F_3$. This applies since $\kappa(H_0)\ge140a\ge10(4a)$.
The unions give the required model and linkage in $F_0$.

\medskip\noindent\emph{Separating and applying induction.}
Fix distinct roots $r'_1,\ldots,r'_a$ in $F_3$. Since
$\delta(F_3)\ge Ka$, choose distinct vertices $s_1,\ldots,s_{2a}$
outside the roots so that $s_j,s_{a+j}$ are adjacent to $r'_j$.
Delete these $3a$ vertices to obtain $F_4$. Apply the separation
hypothesis to $F_4$ and then to one of the two resulting induced
subgraphs, obtaining three disjoint induced subgraphs $J_1,J_2,J_3$.
All these graphs are $K_{14a}$-minor-free. The two applications
are valid because
\[
 \chi(F_4)\ge AT+(B-1025-6K)a,
 \qquad B-1025-6K-14d>28d.
\]
The first separation costs at most $14da$, so the input to the
second has chromatic number at least
$AT+(B-1025-6K-14d)a>28da$. A second loss of at most $14da$
leaves each of the three graphs with
$\chi(J_i)\ge AT+(B-1025-6K-28d)a>7Ka$.
Apply Theorem~\ref{thm:connectivity} inside each $J_i$ with $r=Ka$
and take induced closures to obtain $J'_1,J'_2,J'_3$.

The total chromatic loss for each child is at most
\[
 \underbrace{7a+7a}_{\text{original vertices and proxies}}
 +\underbrace{980a+28a}_{\text{hub and paths}}
 +\underbrace{12Ka}_{\text{connectivity}}
 +\underbrace{3a}_{\text{roots and neighbors}}
 +\underbrace{28da}_{\text{separation}}.
\]
Adding these costs gives
$7+7+980+28+3+28d+12K=1025+28d+12K$.
Consequently
\begin{equation}\label{eq:childbudget}
 \kappa(J'_i)\ge Ka,\qquad
 \chi(J'_i)\ge AT+(B-1025-28d-12K)a
             \ge AT+B(2a/3),
\end{equation}
where the last inequality follows from $B/3>1025+28d+12K$.
The next scale $a'=2a/3$ is integral: if $T=3^m$, then the integer
scales are $(2/3)^iT$ for $0\le i\le m$, and the last is
$2^m\le3^m/\sqrt{m\log3}=T/\sqrt{\log T}$.
This inequality holds for $m\ge5$ (check $m=5$, then use that
$(2/3)^m\sqrt{m\log3}$ decreases). Thus every nonbase scale
occurs before the last integer scale. Induction and~\eqref{eq:childbudget} imply that each
$J'_i$ is $(a',3a')=(a',2a)$-woven.

\medskip\noindent\emph{Assembling the rooted model.}
Choose $a'$ distinct targets in each child, indexed as
$t_1,\ldots,t_{2a}$ in three consecutive blocks of length $a'$.
The graph $F_3-\{r'_1,\ldots,r'_a\}$ has connectivity at least
$(K-1)a\ge20a$. Theorem~\ref{thm:linked} therefore supplies a
linkage from $s_j$ to $t_j$ for $1\le j\le2a$.
Since $|J'_i|\ge Ka+1\ge a'+4a$, Lemma~\ref{lem:simultaneous}
supplies a linkage for the same pairs and rooted $K_{a'}$ models
in the children, meeting the linkage exactly at their targets.

For each $j\in\{1,\ldots,a\}$, unite $r'_j$, its two chosen
neighbor edges, the paths indexed $j$ and $a+j$, and the child
branches containing $t_j$ and $t_{a+j}$. These sets are disjoint
and connected. Each uses two different children, because the two
target indices differ by $a$, exceeding the child block length $2a/3$.
Any two two-element subsets of the three children intersect, and
the corresponding distinct branches in a common child are adjacent.
The resulting sets form a $K_a$ model rooted at the chosen roots.
Thus $F_3$ is $(a,0)$-woven, completing the induction.
\end{proof}

\begin{corollary}\label{cor:outer}
Let $t\ge100$ be an integer, and let $T$ be the least power of three
with $T\ge t$. If a $K_t$-minor-free graph $G$ satisfies the
separation hypothesis of Lemma~\ref{lem:outer} for this $T$ and an
integer $d\ge1$, then
\[
 \chi(G)<3\bigl(10^6(d+1)+62000\bigr)t.
\]
\end{corollary}
\begin{proof}
Let $A,B,K$ be the constants in Lemma~\ref{lem:outer}. If
$\chi(G)\ge(A+B+6K)T$, Theorem~\ref{thm:connectivity} gives an
induced $KT$-connected subgraph $F$ with $\chi(F)\ge(A+B)T$;
its hypothesis holds since $A+B+6K\ge7K$.
Lemma~\ref{lem:outer}, at scale $a=T$, gives a $K_T$ minor,
contradicting $K_t$-minor-freeness. Now use $T<3t$ and $A+6K=62000$.
\end{proof}
\subsection{Completing the proof}

\begin{bootstraptheorem}
If a local linear bound holds at exponent $\alpha>0$, then a local
linear bound holds at exponent $4\alpha/3$.
\end{bootstraptheorem}
\begin{proof}
Let $D,t_0$ be constants for the bound at exponent $\alpha$, and
increase $D$ to an integer.
For sufficiently large $t$, let $G$ be $K_t$-minor-free with
$|G|\le t(\log t)^{4\alpha/3}$, and put
\[
 k=\left\lceil(\log t)^{\alpha/3}\right\rceil\ge2.
\]
Lemma~\ref{lem:packing} gives $W,R,Q$ with
\[
 |Q|\le\frac{|G|}{k}
 \le\frac{t(\log t)^{4\alpha/3}}{(\log t)^{\alpha/3}}
 =t(\log t)^\alpha.
\]
Since $Q$ is $K_t$-minor-free,
the bound at exponent $\alpha$ gives $\chi(G[W])\le2Dt$.

Write $H=G[R]$. We have $b(H)<k$. Let $T$ be the least power of three
at least $t$. At every nonbase scale of Lemma~\ref{lem:outer}, put
$u=14a$. Uniformly throughout these scales,
\[
 \frac{14T}{\sqrt{\log T}}<u\le14T,
 \qquad
 \log u>\log(14T)-\tfrac12\log\log T
 \ge\tfrac12\log t
\]
for all sufficiently large $t$. Moreover,
$k\le2(\log t)^{\alpha/3}$, so
\[
 2k^2\le8(\log t)^{2\alpha/3}
 \le\left(\tfrac12\log t\right)^\alpha
 \le(\log u)^\alpha.
\]
The middle inequality holds eventually because $\alpha/3>0$.
Increasing the lower threshold for $t$ if needed, we have, uniformly
in this window,
\[
 u\ge t_0.
\]
The bound at exponent $\alpha$ therefore colors every $K_u$-minor-free graph
on at most $2uk^2$ vertices with at most $Du$ colors.
Now let $Y$ be an induced $K_u$-minor-free subgraph of $H$ with
$\chi(Y)>28Da=2Du$. Since $b(Y)\le b(H)<k$, we can apply
Lemma~\ref{lem:separation} to $Y$. It gives two disjoint induced
subgraphs with chromatic number at least
$\chi(Y)-Du=\chi(Y)-14Da$. This verifies the separation
hypothesis of Lemma~\ref{lem:outer} with $d=D$ at every nonbase scale.

Corollary~\ref{cor:outer} now gives
\[
 \chi(H)<3\bigl(10^6(D+1)+62000\bigr)t.
\]
Color $W$ and $R$ with disjoint palettes. We obtain
\[
 \chi(G)\le\left[2D+3\bigl(10^6(D+1)+62000\bigr)\right]t.
\]
Both uses of a local coloring bound, on $Q$ and on the small
neighborhood cores, used exponent $\alpha$. This proves
Theorem~\ref{thm:bootstrap}.
\end{proof}

\begin{thebibliography}{99}
\setlength{\itemsep}{2pt}
\bibitem{DP} M.~Delcourt and L.~Postle,
\emph{Reducing linear Hadwiger's conjecture to coloring small graphs},
Journal of the American Mathematical Society 38 (2025), 481--507.
We use the numbering in arXiv:2108.01633v5.
\url{https://arxiv.org/abs/2108.01633v5}.
\bibitem{EGJ} S.~Ehard, S.~Glock, and F.~Joos,
\emph{Pseudorandom hypergraph matchings},
Combinatorics, Probability and Computing 29 (2020), 868--885.
\url{https://doi.org/10.1017/S0963548320000280}.
\bibitem{GN} A.~Gir\~ao and B.~Narayanan,
\emph{Subgraphs of large connectivity and chromatic number},
Bulletin of the London Mathematical Society 54 (2022), 868--875.
\url{https://doi.org/10.1112/blms.12569}.
\bibitem{GK} S.~Gould and T.~Kelly,
\emph{Advancing the R\"odl nibble: new bounds on matchings and the
list chromatic index of hypergraphs},
arXiv:2511.11375v1 (2025).
\url{https://arxiv.org/abs/2511.11375v1}.
\bibitem{Gyarfas} A.~Gy\'arf\'as,
\emph{Problems from the world surrounding perfect graphs},
Zastosowania Matematyki 19 (1987), 413--441.
\bibitem{Had} H.~Hadwiger,
\emph{\"Uber eine Klassifikation der Streckenkomplexe},
Vierteljahrsschrift der Naturforschenden Gesellschaft in Z\"urich
88 (1943), 133--142.
\bibitem{KKMO} D.~Y.~Kang, D.~K\"uhn, A.~Methuku, and D.~Osthus,
\emph{New bounds on the size of nearly perfect matchings in almost
regular hypergraphs},
Journal of the London Mathematical Society (2) 108 (2023), 1701--1746.
\url{https://doi.org/10.1112/jlms.12792}.
\bibitem{Kos} A.~V.~Kostochka,
\emph{Lower bound of the Hadwiger number of graphs by their average degree},
Combinatorica 4 (1984), 307--316.
\url{https://doi.org/10.1007/BF02579141}.
\bibitem{Lin} X.~Lin,
\emph{An asymptotically tight $t\log t$ bound for $k$-connected subgraphs
in dense $K_t$-minor-free graphs},
arXiv:2607.21222v3 (2026).
\url{https://arxiv.org/abs/2607.21222v3}.
\bibitem{LL} C.-H.~Liu and J.~Luo,
\emph{Beyond halfway to Hadwiger's conjecture},
arXiv:2609.06867v2 (2026).
\url{https://arxiv.org/abs/2609.06867v2}.
\bibitem{N} T.~H.~Nguyen,
\emph{Highly connected subgraphs with large chromatic number},
arXiv:2206.00561v2 (2023).
\url{https://arxiv.org/abs/2206.00561v2}.
\bibitem{NPS} S.~Norin, L.~Postle, and Z.-X.~Song,
\emph{Breaking the degeneracy barrier for coloring graphs with no
$K_t$ minor}, Advances in Mathematics 422 (2023), 109020.
\url{https://doi.org/10.1016/j.aim.2023.109020}.
\bibitem{RS} B.~Reed and P.~Seymour,
\emph{Fractional colouring and Hadwiger's conjecture},
Journal of Combinatorial Theory, Series B 74 (1998), 147--152.
\url{https://cgm.cs.mcgill.ca/~breed/SummerNSERC04/frachad.pdf}.
\bibitem{RST} N.~Robertson, P.~Seymour, and R.~Thomas,
\emph{Hadwiger's conjecture for $K_6$-free graphs},
Combinatorica 13 (1993), 279--361.
\url{https://doi.org/10.1007/BF01202354}.
\bibitem{Rodl} V.~R\"odl,
\emph{On a packing and covering problem},
European Journal of Combinatorics 6 (1985), 69--78.
\url{https://doi.org/10.1016/S0195-6698(85)80023-8}.
\bibitem{TW} R.~Thomas and P.~Wollan,
\emph{An improved linear edge bound for graph linkages},
European Journal of Combinatorics 26 (2005), 309--324.
\url{https://doi.org/10.1016/j.ejc.2004.02.013}.
\bibitem{Tho} A.~Thomason,
\emph{An extremal function for contractions of graphs},
Mathematical Proceedings of the Cambridge Philosophical Society
95 (1984), 261--265.
\url{https://doi.org/10.1017/S0305004100061521}.
\bibitem{Wag} K.~Wagner,
\emph{\"Uber eine Eigenschaft der ebenen Komplexe},
Mathematische Annalen 114 (1937), 570--590.
\url{https://doi.org/10.1007/BF01594196}.
\end{thebibliography}
\end{document}
